\section{The problems with industrial agriculture}

\iffalse

LAND USE

Nearly 90 percent of global deforestation is driven by agriculture. Of all agricultural activities, cropland expansion is the single largest contributor, accounting for almost
half of the total deforested area, followed by livestock grazing
-FAO2025StateFoodAgri, page 4

Agricultural expansion is the most widespread form of land-use change, with over one third of the terrestrial land surface being used for cropping or animal husbandry.
-IPBES2019GlobalAssessment, page XVI

Over half the Earth's land surface is under cover types of anthropic origin, including agricultural lands, pasture and cities, leading to loss of natural habitats and displacement of biodiversity.
-IPBES2019GlobalAssessment, page 110

Human drivers extend so widely beyond these areas that as little as 13% of the ocean and 23% of the land is still classified as 'wilderness'.
- IPBES2019GlobalAssessment, page 206

The most accessible and hospitable biomes either have been almost totally modified by humans in most regions (e.g., Mediterranean forests and scrub, temperate forests) or show maximum levels of conversion to anthropogenic biomes or 'anthromes'.
- IPBES2019GlobalAssessment, page 206

Agricultural expansion is one of the main drivers of global biodiversity loss (Eisner et al., 2016; UNCTAD, 2013). In the period 2007-2012, 290,000 km2 of land were cleared for agriculture, a net increase of 29% compared with 2000-2006.
-IPBES2019GlobalAssessment, page 410

Annual tree cover loss increased from 17.2±0.63 million ha/yr in 2001-2010 to 21.3±1.78 million ha/yr in 2011-2016.
-IPBES2019GlobalAssessment, page 410

Industrial agriculture drives land-use change and deforestation, converting natural habitats into cropland and thereby destroying and displacing local biodiversity.
-IPBES2019GlobalAssessment, page 442

Land use change ... resulted in high ecological footprints with increasing pressures on all resources.
-IPBES2019GlobalAssessment, page 503

Land-use and land-cover changes have direct and large impacts on the physical environment.
-IPBES2019GlobalAssessment, page 603

Today, agriculture accounts for 38% of Earth's terrestrial surface (Foley et al., 2011). (look for more recent data)
IPBES2019GlobalAssessment,  page 770

Land-use changes, such as the expansion of soybean in South America, have been identified as a major driver of deforestation and biodiversity loss.
-IPBES2019GlobalAssessment, page 904

By 1990s an area of land the size of China was degraded, and an additional 30% was severely degraded. Over 20% land is now degraded globally. 12 million hectares of land is degraded annually. 50% of irrigated cropland will be salinized by 2050 on current trends. (This is from 2016, check newer data)
-UniformityToDiversity, page 19

Biodiversity loss is the domain in which the world has moved furthest beyond the safe operating space.
-UniformityToDiversity, page 20

Monocropping of feedstocks - accounting for one third of global crops - drives massive biodiversity and habitat loss.
-foodFromSomewhere, page 44

Risk drivers, the root cause of nature loss, are most commonly described based on the IPBES classification (4) - land use and climate change, resource use, pollution and invasive aliens’ species.
-euMethodology, page 8

Agriculture, food and beverage manufacturing, and construction are, for example much more exposed to nature-related risks such as deforestation, than other sub-sectors.
-euMethodology, page 79

Land use change is identified as a key risk driver that contributes to biodiversity loss and ecosystem degradation.
-euMethodology, page 159

In 2023, within agricultural land, world total cropland area was 1 571 million ha (12 percent of total land area). [...] cropland area has grown by 78 million ha since 2001.
The gain in cropland area was largely driven by an increase in permanent crops (such as cocoa, oil palm or coffee), which reached 191 million ha in 2023 (+56 million ha since 2001, a 42 percent increase)
- FAO2025LandStats

glossary:
- NCP: nature’s contributions to people

\fi


\section{The problems of industrial agriculture}
Although the rise of the modern industrial structure of agricultural production has given many societies the ability to forestall famine and thrive in ways long thought impossible, many of its social and ecological repercussions have revealed themselves to be potentially dangerous in the best scenario and significant contributors to an existential crisis in the worst. These problems range from the extensive use of land, to the expulsion of pollutants and greenhouse gasses into the environment, to even social problems like the dependency of a whole population on a few producers. [Give a better overview] The following sections will explore some of these problems.

\subsection{Land Use}

The current paradigm of agricultural production requires the use of vast swathes of land, along with a considerable amount of resources, which inevitably ends up homogenizing the local flora and destroying the habitat of the local fauna, leading to an overall loss of biodiversity. The agricultural expansion that the current system entails is one of the main drivers of global biodiversity loss (\cite[412]{IPBES2019GlobalAssessment}), and is the most widespread form of land-use change, with cropping and animal husbandry covering over one third of the terrestrial land surface (\cite[XVI]{IPBES2019GlobalAssessment}). 

This geographical expansion of the agricultural sector has been a constant for decades, having grown by 78 million ha (780,000 $km^2$) between 2001 and 2023, mostly driven by an increase in permanent crops such as cocoa, oil palm or coffee (\cite{FAO2025LandStats}). All of these are used intensively for monoculture, as are feedstock crops which by themselves account for one third of the total crops produced globally, this in turn drives massive biodiversity and habitat loss (\cite[22]{foodFromSomewhere}). Monoculture is another significant problem that will be explored in a subsequent section.

According the the FAO's State of Food and Agriculture report (\cite[4]{FAO2025StateFoodAgri}), agriculture currently drives 90 percent of global deforestation, with cropland expansion accounting for almost half of the total deforested area.  [Find source to quantify deforestation impact on climate change]













